\documentclass[11pt,a4paper]{article}
\usepackage[margin=25mm]{geometry}
\usepackage{amsmath,amssymb,amsthm,booktabs}
\newtheorem{lemma}{Lemma}\newtheorem{theorem}{Theorem}
\title{Odd kernels and two pair counts with an element of order seven}
\author{Human Compression project}\date{}
\begin{document}\maketitle
Let $G=O'N$, and let $x_J(p,q)$ and $e_J(p,q)$ count respectively
all ordered pairs $(r,s)$ in $J$ of orders $p,q$ with $|rs|=2$, and
those which generate $J$. We recover the fixed Connor--Leemans
Table~1 entries
\[
                 e_G(4,7)/(2|G|)=102,\qquad
                 e_G(7,7)/(2|G|)=648.                         \tag{1}
\]
The complete subgroup roster, primary ATLAS data, accepted small
group models and prior exact character evaluations are explicit
borrowed premises. Row numbers refer to the fixed subgroup roster.

\begin{lemma}\label{zero}
The point and plane parabolics in rows 9 and 10, and the flag
normalizer in row 22, have no pair of type $(4,7)$ or $(7,7)$.
\end{lemma}
\begin{proof}
The accepted projective identification gives a normal subgroup
$V\cong C_7^2$ in a point or plane parabolic $P$, with quotient
$P/V\cong\mathrm{GL}_2(7)/C_3$. Projection through the odd kernel
preserves the exact order of every element of 2-power order, since
its cyclic subgroup meets the kernel trivially. In particular, an
involutory product remains a nonidentity involution, and an element
of order four retains order four. An element of order seven has
image of order one or seven.

Put $Q=\mathrm{GL}_2(7)/C_3$. Its determinant map has image of order
two after the scalar quotient, and its kernel is the embedded
$D=\mathrm{SL}_2(7)$. Every element of order seven belongs to $D$.
Every element of order four does also: choose a matrix lift $A$;
its fourth power is a scalar in $C_3$, which can be removed by
multiplying by a unique scalar in $C_3$, because fourth powering
is a bijection there. The resulting lift has order four. Over
$\mathbb F_7$, its primitive fourth-root eigenvalues form the
conjugate pair $i,-i$, so its determinant is one. A split or scalar
matrix cannot have order four because $\mathbb F_7^*$ has order six.
The determinant-one group $D$ has exactly one involution, $-I$:
an involution is diagonalizable, and determinant one forces either
$I$ or $-I$.

If the two projected factors both have order seven, their product
cannot be $-I$. Indeed this would give $\bar s=-\bar r^{-1}$,
of order fourteen. If one projected factor is the identity, its
product with a factor of order one or seven cannot be an involution.
For type $(4,7)$, the equation $\bar r\bar s=-I$ would instead
give $\bar s=-\bar r^{-1}$ of order four, again impossible.
This excludes both types in $P$.

The flag normalizer has a normal Sylow 7-subgroup of order $7^3$,
as its upper triangular projective model shows. Two order-seven
factors have product in that odd subgroup. For type $(4,7)$,
projection modulo it preserves order four and kills the order-seven
factor, so the product cannot have order two. This proves the lemma.
\end{proof}

\begin{lemma}\label{atmost}
A pair of type $(4,7)$ or $(7,7)$ has at most one common commuting
involution. It therefore lies in at most one maximal subgroup of row 7.
\end{lemma}
\begin{proof}
The centralizer of an order-seven element has order 49 or 1372.
In the latter case the primary power data supply an order-28 root,
so a Sylow 2-subgroup of the centralizer is cyclic of order four.
Suppose that the common centralizer of the pair had an element $v$
of order seven. Both pair entries would then lie in $C_G(v)$.
This group has a normal Sylow 7-subgroup: the number of such Sylow
subgroups divides four and is congruent to one modulo seven.
Two order-seven entries cannot have an involutory product there.
For type $(4,7)$, projection modulo that odd subgroup preserves the
order-four entry and kills the other entry, again contradicting an
involutory product. Thus the common centralizer has no element of
order seven. Its order divides 49 or 1372, so it is a 2-group and
is contained in a cyclic Sylow 2-subgroup. It has at most one
involution. Row 7 consists of the involution centralizers.
\end{proof}

\begin{lemma}\label{inside}
If such a pair lies in a projective maximal $M$ of row 2 or 3 and
in $C_G(t)$ of row 7, then $t\in M$. No pair involving an order-seven
entry lies both in a row 7 maximal and in one of rows 6, 21, 43 or 44.
\end{lemma}
\begin{proof}
The second assertion is the empty intersection of the cyclic-seven
ancestor set with the descendants of row 7 and each listed maximal.
For the first, intersect the two relevant cyclic ancestor sets with
the descendants of rows 7 and 2 or 3, and remove the descendants of
the flag normalizer, which Lemma~\ref{zero} excludes. For both types
the possible generated-subgroup classes are
\[
65,66,110,111,112,113,115,116,117,118,171,173,175,176,177,242.
\]
For every class except 177 and 242, reuse the accepted common-involution
argument for type $(7,8)$. In detail, a subgroup of index at most four
in the block group of order 1344 contains its perfect subgroup
$\mathrm{SL}_2(7)$ and hence its central involution. This treats
65, 66, 112, 113, 115--118, 175 and 176. The groups 110 and 111 are
$\mathrm{PGL}_2(7)\times C_2$, with a central involution. Every
projective maximal containing a group of class 171 or 173 contains
its unique parent of class 110 or 111: that parent is its full
normalizer and normalizes its unique simple overgroup. These
arguments depend only on the class and on uniqueness of the common
involution, which Lemma~\ref{atmost} supplies here.

Class 177 is $\mathrm{PSL}_2(7)\times C_2$, and its central involution
belongs to every overgroup containing it. This structure can be
read within its parent $\mathrm{PGL}_2(7)\times C_2$: among the
index-two subgroups, it is the one containing the central involution
and having derived subgroup of order 168. The other two are the
centerless classes 171 or 173. The complete maximal-child roster
identifies the remaining class as 177.
A subgroup of class 242 has immediate overgroups only in 171, 173
and 177. Along the first two routes, the preceding normalizer
argument puts the appropriate centralizer group 110 or 111 inside
the projective maximal; along the third, that maximal contains 177.
Each route therefore puts an involution centralizing the generated
subgroup inside $M$. By Lemma~\ref{atmost} this involution is $t$.
\end{proof}

Write $L\cong\mathrm{PSL}_2(7)$ and $F\cong C_2^3:C_7$, with the
nontrivial irreducible action of $C_7$ on $C_2^3$.
\begin{lemma}\label{outside}
Among proper subgroups generated by a pair of type $(4,7)$ and
contained in no row 7 maximal, only classes 239, 240 and 241 have
more than one maximal overgroup. They are copies of $L$, and their
numbers of actual maximal overgroups are respectively three, two
and three. For type $(7,7)$ there is one further possibility, class
374, a copy of $F$, with exactly two maximal overgroups.
\end{lemma}
\begin{proof}
Use the cyclic ancestor join, removing $G$, descendants of row 7,
and descendants of rows 9, 10 and 22 by Lemma~\ref{zero}. For type
$(4,7)$ this leaves
\[
2,3,4,5,21,43,44,68,145,172,174,239,240,241.
\]
The row 21 maximal has zero character count for this type, excluding
it and its descendants 68 and 145. Classes 4, 5, 172 and 174 have
unique actual maximal overgroups, by the previously checked direct
class-length incidences. The others, apart from the three stated
exceptions, are themselves maximal.

For type $(7,7)$ the corresponding set is
\[
2,3,4,5,6,21,43,44,68,145,172,174,237,239,240,241,374.
\]
We must exclude full generation of classes 68, 145 and 237.
In the accepted row 21 model, the normal kernel is $A\cong C_4^3$
and the quotient is $L_3(2)$. The preimage of the normalizer $7:3$
of a Sylow 7-subgroup of that quotient is a maximal subgroup of
order 1344; the complete row 21 maximal list identifies it as row 68.
Its preimage of the normal subgroup of order seven has order 448
and is row 145. Thus 68 has a quotient of order three, excluding
generation by two elements of order seven. A row 145 group is
$A:C_7$: its complement exists by coprime splitting. If two
order-seven elements $r,s$ in it have product $w$ of order two,
then $w$ lies in $A$, and hence in $A[2]\cong C_2^3$.
Their generated group is $\langle r,w\rangle$, contained in
$A[2]\langle r\rangle$, of order at most 56. It cannot be the
group of order 448. A row 237 group has its unique subgroup of
order 56, class 374, as a normal subgroup of index three, so the
same quotient argument excludes generation of row 237.

For class 239, direct incidence counts give two actual A7 overgroups
in row 43 and one immediate parent in row 174, whose unique maximal
parent is row 3. Thus there are three maximal overgroups.
Class 241 has the symmetric list, with two row 44 overgroups and
one row 2 overgroup. Class 240 has just one overgroup of each A7
class. The complete parent lists exclude other routes. The three
class lengths are $|G|/336$, $|G|/168$ and $|G|/336$.

Finally, a row 374 group $F_0$ has unique actual immediate parents
$D_0$ and $K_0$ in rows 145 and 237. The subgroup $F_0$ is normal
in $K_0$, since the latter has exactly one maximal subgroup of
order 56. As $|N_G(F_0)|=168=|K_0|$, we have $K_0=N_G(F_0)$.
Uniqueness of $D_0$ forces $K_0$ to normalize it. The normalizer
of $D_0$ has order 1344 and is its unique row 68 parent. Hence the
two routes from $F_0$ through 145 and 237 meet in the same row 68
subgroup and lead to its unique row 21 maximal. The row 237 subgroup
also has a unique row 6 maximal. There are no further routes.
Consequently $F_0$ has exactly two maximal overgroups, and its
class length is $|G|/168$.

For clarity, the isomorphism type used here follows from the actual
subgroup roster, not merely its abbreviated label. A row 374 group
has exactly one maximal subgroup of order eight, which has no
cyclic subgroup of order four and is therefore elementary abelian,
and eight maximal cyclic subgroups of order seven. Thus the order-eight
subgroup is normal, the split action of a complement is nontrivial,
and the action on its three-dimensional $\mathbb F_2$ space is
irreducible. All subgroups of order seven in $\mathrm{GL}_3(2)$ are
conjugate by Sylow's theorem. This identifies $F_0$ with $F$.
\end{proof}

\begin{theorem}
Let $H$ and $K$ be the accepted inner and outer involution centralizers
in $L_3(7):2$, of orders 1344 and 672. For $p\in\{4,7\}$,
\begin{align*}
 e_G(p,7)={}&x_G(p,7)-\sum_M[G:M]x_M(p,7)\\
 &+2\cdot342868680\bigl(x_H(p,7)+2x_K(p,7)\bigr)\\
 &+\frac{|G|}{56}e_L(p,7)
       +\begin{cases}0,&p=4,\\ |G|e_F(7,7)/168,&p=7.\end{cases}       \tag{2}
\end{align*}
The resulting ordered automorphism-orbit counts are those in (1).
\end{theorem}
\begin{proof}
For a pair having a row 7 overgroup, Lemmas~\ref{atmost} and
\ref{inside} reduce the excess in the maximal sum to the projective
centralizer contributions. As in the accepted $(7,8)$ formula,
each of the two projective classes has 122760 maximal subgroups,
each with 2793 inner involutions and 5586 outer involutions.
Their product gives the second line of (2).

For pairs with no row 7 overgroup, Lemma~\ref{outside} lists every
generated subgroup having excess multiplicity. A pair determines
its generated subgroup uniquely. Thus different actual copies of
$L$ or $F$ supply disjoint contributions. The three $L$ classes
contribute
\[
 \left(2\frac{|G|}{336}+\frac{|G|}{168}
                  +2\frac{|G|}{336}\right)e_L(p,7)
                       =\frac{|G|}{56}e_L(p,7).
\]
The additional $F$ class contributes $|G|e_F(7,7)/168$.

The literal certificates give
\[
\begin{array}{c|rrrr}
p&x_H(p,7)&x_K(p,7)&e_L(p,7)&e_F(p,7)\\\hline
4&0&672&336&0\\
7&0&336&336&336.
\end{array}
\]
Here all counted pairs in $L$ generate it: its maximal subgroups
are $S_4$ and $7:3$, each lacking one necessary element order.
Every counted $(7,7)$ pair in $F$ generates it because its proper
maximal subgroups have orders eight and seven. The exact character
evaluation of the first line of (2) and these corrections give (1)
after division by $|\operatorname{Aut}(G)|=2|G|$. This is the free
action on ordered generating pairs, with no quotient by duality.
\end{proof}

The finite reader reconstructs the two retained centralizers, the
square-determinant subgroup of the retained projective line action,
and the 56 affine permutations of $\mathbb F_8$. It checks all
order-seven classes and every partner of each representative.
The complete flag model independently confirms the zero counts
in Lemma~\ref{zero}. The present contribution is the odd-kernel
exclusion and the exact correction by uniquely generated subgroups.
The numerical assertions themselves are due to Connor and Leemans.
No new chirality, whole-source, open-question or Lean credit is claimed.
\end{document}
